\scanpage{117}
\mspage{111}
% archive sheet number: 116

Ainsi, le premier axiome sur les supports dans un atelier, muni d'un magasin \deleted{\unsure{adjoint}}\inserted{associé} $\unsure{\mathscr{F}|\mathcal{F}|F}$, sera

\uline{Axiome supp 1} Pour toute figure $F\in\mathscr{F}_{c}$\inserted{$\operatorname{Figat}(\unsure{\mathcal{A}|\mathscr{A}|\mathbb{A}})$}, la famille des $(X^{\unsure{\circ|0|o}})_{X\in F}$ de $\Sigma_{\mathcal{A}}$ est admissible.

Cet axiome n'est pas \unsure{redondant}; comme on voit en \unsure{partant} d'un \unsure{qtopos|topos} $\Sigma$, et en prenant l'atelier $\operatorname{Figel}\overset{\deleted{\mathrm{fin}}}{\operatorname{spat}}_{\mathrm{fin}}(\Sigma)$, \deleted{et} des figures élémentaires\inserted{spatiales} finies (disons) de $\Sigma$, et le magasin \deleted{des} associé $\operatorname{Fig}\deleted{\mathrm{el}}\operatorname{spat}_{\mathrm{fin}}(\Sigma)$ des figures spatiales finies. La validité de l'axiome des supports revient alors, d'après le cor. 1, à ceci
\[
\left\vert\begin{aligned}
&\text{si }e,e'\in\Sigma\text{ avec }e\mid\unsure{0|o}\mid e',\quad
\deleted{\text{et si }e''=\unsure{\complement|C}e\vee e',}\\
&\deleted{\text{alors}}\ \text{alors }e'=(\complement e)\wedge(e\vee e').
\end{aligned}\right.
\]
Ou encore, la validité des conditions du lemme 2 (p. 105) sur $\Sigma$.
